\section{Cơ sở mô hình}
\label{sec:background}

\subsection{BiLSTM}

BiLSTM xử lý chuỗi theo cả hai chiều thời gian, nhờ đó mỗi biểu diễn có thể tiếp nhận ngữ cảnh trước và sau vị trí đang xét. Đây là mô hình chuỗi được sử dụng trong cấu hình đối chứng và là mốc để đánh giá khả năng thay thế bằng mạng tích chập theo thời gian.

\subsection{Mạng tích chập theo thời gian}

Mạng tích chập theo thời gian sử dụng các phép tích chập giãn và cấu trúc nhân quả để mở rộng trường tiếp nhận theo chuỗi. So với mô hình hồi tiếp, cách biểu diễn này có tiềm năng xử lý song song và giảm độ trễ suy luận. Tuy nhiên, lợi thế kiến trúc chỉ có ý nghĩa khi được xác nhận bằng so sánh bắt cặp trên cùng dữ liệu.

\subsection{ResNet-1D}

ResNet-1D sử dụng các kết nối dư để hỗ trợ tối ưu hóa mạng tích chập sâu trên tín hiệu một chiều. Trong nghiên cứu, ResNet-1D được dùng làm bộ trích đặc trưng trước mô hình chuỗi. Hiệu quả của nó được đánh giá ở hai mặt: chất lượng dự đoán và khả năng tạo ra biểu diễn hữu ích cho nhiệm vụ phân loại chuỗi.

\subsection{Chỉ số đánh giá}

Macro-F1 là trung bình F1 của năm lớp và được chọn làm chỉ số chính vì phản ánh cân bằng giữa các lớp không đồng đều. Accuracy và Cohen's kappa được báo cáo bổ sung. F1 theo lớp, đặc biệt là N1, được dùng để mô tả các thay đổi mà chỉ số tổng hợp có thể che khuất.
